\begin{abstract}
Let a directed cycle of length (n\geq 3) be visible, and place a binary
receipt fiber over each visible state. We classify every deterministic lift
of the visible successor that is covariant under the common binary deck
transformation. Such a lift is specified by one binary voltage on each
visible edge. Fiber-preserving changes of local sheet labels act as gauge
transformations, and the sum of the edge voltages around the circuit is a
complete gauge invariant. Hence exactly two gauge classes occur. Trivial
circuit holonomy produces two disjoint directed (n)-cycles. Nontrivial
circuit holonomy produces one directed (2n)-cycle, exchanges the receipt
sheet after one visible circuit, and returns to the initial lifted state only
after two visible circuits. Therefore a visible order-(n) cycle together
with a nontrivial binary full-return receipt uniquely forces a connected
order-(2n) lift, up to fiber-preserving gauge and isomorphism. The familiar
(C_4\) to (C_8\) lift is the case (n=4), not an exceptional coincidence.
The theorem is purely finite and algebraic; no physical realization is
claimed.
\end{abstract}

